\documentclass[11pt]{article}
\usepackage{mathpartir}
\usepackage{pifont}
\usepackage{multicol}
\usepackage{listings}
\usepackage{pgf}
\usepackage{tikz}
\usepackage{alltt}
\usepackage{hyperref}
\usepackage{url}
\usepackage{amsmath}
\usepackage{amssymb}
\usepackage{dafny}
\usetikzlibrary{arrows,automata,shapes,positioning}
\tikzstyle{block} = [rectangle, draw, fill=blue!20, 
    text width=5em, text centered, rounded corners, minimum height=2em]
\tikzstyle{bt} = [rectangle, draw, fill=blue!20, 
    text width=1em, text centered, rounded corners, minimum height=2em]
\newcommand{\xmark}{\ding{55}}

\newtheorem{defn}{Definition}
\newtheorem{crit}{Criterion}
\newcommand{\true}{\mbox{\sf true}}
\newcommand{\false}{\mbox{\sf false}}

\newcommand*\circled[1]{\tikz[baseline=(char.base)]{
            \node[shape=circle,draw,inner sep=2pt] (char) {#1};}}


\newcommand{\handout}[5]{
  \noindent
  \begin{center}
  \framebox{
    \vbox{
      \hbox to 5.78in { {\bf Software Testing, Quality Assurance and Maintenance } \hfill #2 }
      \vspace{4mm}
      \hbox to 5.78in { {\Large \hfill #5  \hfill} }
      \vspace{2mm}
      \hbox to 5.78in { {\em #3 \hfill #4} }
    }
  }
  \end{center}
  \vspace*{4mm}
}

\newcommand{\lecture}[4]{\handout{#1}{#2}{#3}{#4}{Lecture #1}}
\topmargin 0pt
\advance \topmargin by -\headheight
\advance \topmargin by -\headsep
\textheight 8.9in
\oddsidemargin 0pt
\evensidemargin \oddsidemargin
\marginparwidth 0.5in
\textwidth 6.5in

\parindent 0in
\parskip 1.5ex
%\renewcommand{\baselinestretch}{1.25}

\usepackage{enumitem}

\newtheorem{prop}{Proposition}
\newtheorem{lemma}{Lemma}
\usepackage{ebproof}
\newcommand{\qedsymbol}{\rule{1ex}{1ex}}
\newcommand{\sem}[3]{\langle #1, #2 \rangle \Downarrow #3}

\lstset{ %
language=Java,
basicstyle=\ttfamily,commentstyle=\scriptsize\itshape,showstringspaces=false,breaklines=true,numbers=left}

%\usepackage{fontspec}
%\setmonofont{Cousine}[Scale=MatchLowercase]

\begin{document}

\lecture{20 --- March 20, 2026}{Winter 2026}{Patrick Lam}{version 1}

\section*{Termination}
Dafny does prove termination for loops and recursion. It does so by proving that there is something
that decreases and that is bounded. Usually Dafny manages to guess what is decreasing. When that is
something over natural numbers, the lower bound is often 0. You can hover over
the loop to see the inferred decreases clause. Sometimes Dafny can't guess, and you have to provide an annotation.

\begin{lstlisting}[language=dafny]
method Decreases()
{
  var i := 20;
  while 0 < i
    invariant 0 <= i
    decreases i // could be inferred
    {
      i := i - 1;
    }
}
\end{lstlisting}
Note that the thing that is decreasing might be the negation of something that's increasing.
\begin{lstlisting}[language=dafny]
method Decreases2()
{
  var i, n := 0, 20;
  while i < n
    invariant 0 <= i <= n
    decreases n - i
    {
      i := i + 1;
    }
}
\end{lstlisting}
Dafny can guess this one too, from the loop condition $i < n$.
The lower bound is from $i \leq n$ in the invariant, which implies
$n - i \geq 0$. The bound doesn't have to be constant, as in binary search.

\paragraph{Exercise 10.} In the loop above, the invariant $i \leq n$
and the negation of the loop guard allow us to conclude $i == n$ after the loop,
as we've previously checked with an \textsf{assert}. If the loop guard was instead
$i \neq n$ (as in Exercise 7), then the negation of the guard immediately gives $i == n$,
without needing the loop invariant. Change the loop guard to $i \neq n$ and delete the invariant.
What happens?

We have also seen the recursive function \textsf{fib}. Dafny also guesses the decreases
annotation here. We could specify it explicitly:

\begin{lstlisting}[language=dafny]
function fib(n:nat): nat
  decreases n
{
  if n == 0 then 0
  else if n == 1 then 1
  else fib(n-1) + fib(n-2)
}
\end{lstlisting}

There is a longer Dafny tutorial on termination, and we'll probably go through that one as well.

\section*{Arrays}
In Dafny we have one-dimensional arrays $\mathsf{array}\langle T \rangle$ for any $T$ (though here we'll only consider
$\mathsf{array}\langle int \rangle$), as well as $\mathsf{array}?\langle T \rangle$ which includes possibly-null arrays.

Array access is as you might expect, $a[i]$, and there is a built-in
length field: $\mathsf{a}.\mathsf{Length}$.  Part of Dafny's
compile-time verification includes showing that all array accesses are
in-bounds, i.e.  betweeen 0 and $\textsf{Length}-1$ inclusive.

We saw array allocation right at the beginning of this document:
\begin{lstlisting}[language=dafny]
    var arr := new int[2];
\end{lstlisting}

Let's express the easy part of the contract for the \textsf{Find} method.
\begin{lstlisting}[language=dafny]
method FindPartialContract(a: array<int>, key:int) returns (index:int)
  // partial contract
  ensures 0 <= index ==> index < a.Length && a[index] == key
{
  // can you write code that satisfies this postcondition?
  // it can be done with one statement.
}
\end{lstlisting}
This is the case where the key is in the array. It's at index \textsf{index}.
The conjunct $\mathsf{index} < a.\mathsf{Length}$ guards against out-of-bounds array
accesses that might be caused by \textsf{a[index]}, and works because Dafny has short-circuit
logical ands/ors (like C and many other languages).

For the other case, we need to say that \textsf{key} does not exist at \emph{any}
index in the array, so we need quantifiers.

\paragraph{Quantifiers.} Dafny supports quantifiers in expressions. Use a double-colon (::)
after the variable name; it shows up as a dark $\cdot$ in Emacs and VS Code.

\begin{lstlisting}[language=dafny]
method WithQuantifier()
{
  assert forall k :: k < k + 1;
}
\end{lstlisting}
(This verifies, but there's a warning about triggers. We won't talk about this here.)

Typically we quantify over some more limited set than ``all integers'' (Dafny infers the type of
$k$, though you can specify it if you want). For instance, all integers that are valid indices into an array,
i.e. between 0 and $a.\mathsf{Length}$.
\begin{lstlisting}[language=dafny]
  forall k :: 0 <= k <= a.Length => a[k] != key
\end{lstlisting}
We can thus write the whole contract:
\begin{lstlisting}[language=dafny]
method Find(a: array<int>, key:int) returns (index:int)
  ensures 0 <= index ==> index < a.Length && a[index] == key
  ensures index < 0 ==> forall k :: 0 <= k < a.Length ==> a[k] != key
{
    // TODO implement this
}
\end{lstlisting}
Linear search is the easiest way to implement \textsf{Find}:
\begin{lstlisting}[language=dafny]
  index := 0;
  while index < a.Length {
    if a[index] == key {
      return;
    }
    index := index + 1;
  }
  index := -1;
\end{lstlisting}
and when we plug that into the contract, we see that Dafny can't prove the postcondition
saying that if \textsf{index} is negative, then \textsf{key} is not in the array. Makes sense:
all we know upon leaving the loop is that \textsf{index} has reached the end of the array.
We need to add an invariant showing that the part of the array we've visited so far does not
contain \textsf{key}.
\begin{lstlisting}[language=dafny]
  invariant forall k :: 0 <= k < index ==> a[k] != key
\end{lstlisting}
Dafny complains: it doesn't know that \textsf{index} doesn't go off the end of the array.
This invariant works.
\begin{lstlisting}[language=dafny]
  invariant 0 <= index <= a.Length
\end{lstlisting}
We often see loop invariants where we build the property that we need one element at a time,
as we do here. In this case, the property is that the visited part of the array doesn't contain
the key that we're looking for.

\paragraph{Exercise 11.} Write a method that takes an integer array, which it requires to have
at least one element, and returns an index to the maximum of the array's elements. Annotate
the method with pre- and postconditions that state the intent of the method, and annotate
its body with a loop invariant to verify it.
\begin{lstlisting}[language=dafny]
method FindMax(a: array<int>) returns (index:int)
  // TODO put preconditions and postconditions here  
{
  // TODO implement this
}
\end{lstlisting}

Linear search is fine, but if the array is sorted, then binary search is faster. We can't
prove that binary search is correct unless we can somehow express a property about the array
being sorted. Which we can. But let's see if we can refactor this property, using predicates.

\paragraph{Predicates.} Dafny has functions, which can return anything. Dafny also has
predicates, which return booleans. We are going to define the \textsf{sorted} predicate
which will return true if and only iff its array argument is sorted. Using this predicate
makes our expressions shorter and more readable (as long as the name of the predicate
corresponds to its meaning).

\begin{lstlisting}[language=dafny]
predicate sorted(a: array<int>)
{
  forall j, k :: 0 <= j < k < a.Length ==> a[j] <= a[k]
}
\end{lstlisting}
A predicate is like a function, but the return type must be boolean, and so you don't have to write this.

Dafny won't compile this, and says that there is an insufficient \textsf{reads} clause.

\paragraph{Framing.} Dafny needs to know when heap locations might change.

We know that functions and predicates can't change the heap. But
Dafny isn't a purely functional language---methods can change the heap.

So, Dafny wants to know when the value of \textsf{sorted} might change
between calls. In a world where there are two arrays, and a method
writes to one of the arrays, a \textsf{reads} annotation tells
Dafny \emph{which} array \textsf{sorted} is reading from, and whether
that array might have been changed by the method. If not, then the
array has the same sortedness before and after the method was called.

\begin{lstlisting}[language=dafny]
predicate sorted(a: array<int>)
  reads a
{
  forall j, k :: 0 <= j < k < a.Length ==> a[j] <= a[k]
}
\end{lstlisting}
The \textsf{reads} annotation specifies a set of memory
locations that the function is allowed to access. Here, it is
an array, which means that \textsf{sorted} can access all
of the elements of the array. It could also be object fields and
sets of objects. We'll probably check that later.

As we've seen, Dafny checks that the \textsf{reads} clause
specifies everything that the function might read. This includes
any functions that the function transitively calls.

Conversely, as we've seen, methods can read anything they want.
But they must declare any changes to memory using a \textsf{modifies}
clause.

These clauses allow Dafny to reason about methods and functions
in isolation, and hence they make verification much more tractable.

Dafny only requires framing information for memory on the heap,
which includes arrays and objects, but not sets, sequences, or multisets
(which we'll talk about later).

\paragraph{Exercise 12.} Create a \textsf{sortedAndDistinct}
predicate that returns true exactly when the array is sorted and all its
elements are distinct.

~\\[-5em]
\paragraph{Exercise 13.} Create a \textsf{sortedAndNonNull}
predicate that accepts references to a possibly-null array and that
returns true exactly when the array is sorted and not null.

\paragraph{Binary Search.} Predicates: software engineering principles for
annotations. Here we use the \textsf{sorted} predicate in the precondition
for \textsf{BinarySearch}. Dafny thus knows that the array is sorted
when verifying the method body.

\begin{lstlisting}[language=dafny]
method BinarySearch(a: array<int>, value: int) returns (index:int)
  requires 0 <= a.Length && sorted(a)
  ensures 0 <= index ==> index < a.Length && a[index] == value
  ensures index < 0 ==> forall k :: 0 <= k < a.Length ==> a[k] != value {
    // ...
}
\end{lstlisting}

Here's an implementation of binary search.
\begin{lstlisting}[language=dafny]
  var low, high := 0, a.Length;
  while low < high
    invariant 0 <= low <= high <= a.Length
    invariant forall i ::
      0 <= i < a.Length && !(low <= i < high) ==> a[i] != value
    {   
      var mid := (low + high) / 2;
      if a[mid] < value {
        low := mid + 1;
      } else if value < a[mid] {
        high := mid;
      } else {
        return mid;
      }
    }
    return -1;
\end{lstlisting}
Let's look at the implementation and the invariants.
\begin{lstlisting}[language=dafny]
    invariant 0 <= low <= high <= a.Length
    invariant forall i ::
      0 <= i < a.Length && !(low <= i < high) ==> a[i] != value
\end{lstlisting}
The first invariant says that the search is looking between
$[\mathsf{low}, \mathsf{high})$, and the second invariant
says that the searched-for value is not anywhere outside
that range.

Reasoning about how the first invariant continues to hold is
straightforward.

For the second invariant, Dafny needs to show that changing
\textsf{low} or \textsf{high} doesn't miss the value that we are
searching for. It can use the fact that predicate \textsf{sorted}
holds. Taking the first \textsf{if} case, we are bumping up
\textsf{low} to above \textsf{mid}, based on sortedness and the fact
that $\mathsf{a}[\mathsf{mid}] < \mathsf{value}$, which ensure that
the value isn't in that part of the range.

Off-by-one errors are easy to write. Dafny catches them in this case.

~\\[-4em]
\paragraph{Exercise 14.} Change the assignments in the body of
\textsf{BinarySearch} to set \textsf{low} to \textsf{mid} or
to set \textsf{high} to $\mathsf{mid}-1$. What goes wrong?

\section*{Conclusion}
We've seen a bunch of the major features of Dafny and used it to verify
some simple code, up to a binary search implementation. There are more
features: objects, sequences and sets, data structures, lemmas, modules, etc.

There is a lot of information about Dafny online, including other
tutorials and the reference.

SE 465 is not actually the course on specifications. That's SE 463.
But we use formal specifications in this part of SE 465 to prove
that implementations conform to these specifications, and that
they are using other code consistently with that code's specifications.
Both tests and verification are tools that you can use to verify
that your code does what it's supposed to do. Tests are more
practical, while verification is conceptually more powerful.

\end{document}
